\section{Mechanism Analysis}\label{sec:theory}
\subsection{Latent Scores and Binary Contrast}
Assume fixed dynamics, a frozen prior, an exogenous offset law, and regularity allowing differentiation under the expectation. The augmented trajectory density then factors into exogenous terms and conditional latent-policy terms. Consequently,
\begin{equation}
\nabla J=\mathbb E\!\left[Y\sum_t\nabla\log\pi_\theta(u_t\mid h_t,p_t,b_t)\right].
\label{eq:score}
\end{equation}
To see this, differentiate the trajectory density: only the policy factors depend on $\theta$ at fixed recorded history, while clipping is a deterministic execution map. The offset density cancels in likelihood ratios. A learned covariance contributes its own Gaussian score. Equation~\eqref{eq:score} describes the full-trajectory score objective, not an unbiasedness guarantee for the stage-filtered PPO surrogate~\cite{capg}.

For $G$ conditionally independent candidates of one scene with binary milestone probability $p$, the probability of a nonconstant group is
\begin{equation}
P_{\rm mixed}=1-p^G-(1-p)^G.
\label{eq:mixed}
\end{equation}
This follows by excluding the all-success and all-failure outcomes. At $G=8$, changing $p$ from 0.03 to 0.30 increases this probability from 21.6\% to 94.2\%. Thus milestones can supply contrast when completion is rare, although an almost certain milestone also loses contrast. For heterogeneous scenes, average the scene-specific probabilities rather than inserting the overall success rate. These observations agree with binary-group analyses~\cite{mroueh2026,bay2026}. They connect the studies: milestones distinguish different achievement levels, while dense progress can distinguish trajectories within one level. Neither guarantees useful updates: progress must vary, and rollout filtering must retain the group. In particular, a terminal-accuracy filter can discard an all-failure group even when its dense returns differ.

\subsection{What Prefix Protection Guarantees}
Consider a benchmark of three sequential Bernoulli gates with separate logits, $q_k=\operatorname{sigmoid}(\theta_k)$ and $J=q_1q_2q_3$. A reached gate produces one decision $X_k$, with score $X_k-q_k$. Use sign-preserving binary group normalization, nonnegative weights, on-policy raw ascent, $\beta\ge0$, and $\kappa=0$, without regularization or optimizer momentum. Then every realized batch direction has $d_k\ge0$, and
\begin{equation}
J(\theta+\eta d)\ge J(\theta),\qquad\eta\ge0.
\end{equation}
Indeed, an achieved gate receives nonnegative milestone credit and no negative completion bonus; a failed gate receives nonpositive credit. In either case credit and score have the same sign. Their weighted product is nonnegative, so every gate logit can only increase. This is a sufficient explanation of the protection rule, not evidence that symmetric credit must fail in this model.

The neural controller shares parameters across stages and assigns many actions to each stage, so the sign argument no longer applies. For an objective with an $L$-Lipschitz gradient on the actual update segment,
\begin{equation}
J(\theta+\Delta)-J(\theta)\ge\nabla J^\top\Delta-\tfrac L2\|\Delta\|^2.
\end{equation}
The inequality follows by integrating the gradient along that segment. Positive alignment and a sufficiently small update would imply improvement, but neither an alignment bound nor $L$ has been measured for the CDPR policy. PPO clipping alone does not establish this condition~\cite{trpo}. The analysis also gives no automatic guarantee for hardware transfer or for deployment with a different exploration law.
